\begin{remark}\label{sga3:XXI:7.2.4}
Reasoning similarly gives the following result.  If \(\cR\) is
irreducible and nonreduced, with \(\rgss(\cR)>1\), then
\[
 \operatorname{long}(R)=\{1,2,4\}.
\]
If \(R_i=\operatorname{long}^{-1}(i)\), then
\[
 \operatorname{ind}(R)=R_1\cup R_2,\qquad R_4=2R_1,
\]
and two nonproportional roots of \(R_1\) are orthogonal.
Conversely, suppose that \(R\) is an irreducible reduced root
system such that \(\operatorname{long}(R)=\{1,2\}\).  Put
\(R_i=\operatorname{long}^{-1}(i)\), and suppose that two
nonproportional roots of \(R_1\) are orthogonal.  Then
\(R\cup2R_1\) is irreducible and nonreduced, and
\[
 \operatorname{ind}(R\cup2R_1)=R.
\]
\end{remark}

\begin{lemma}\label{sga3:XXI:7.2.5}
If \(\cR\) is an irreducible root datum, then \(\cR^\ast\) is
irreducible as well, and the product
\[
 \operatorname{long}(\alpha)\operatorname{long}(\alpha^\ast)
\]
is constant as \(\alpha\) runs through \(R\).
\end{lemma}

\begin{proof}
This follows immediately from \ref{sga3:XXI:7.1.3}\,\SGARefLink{sga3:XXI:7.2:item:iv}{(iv)} and
\ref{sga3:XXI:2.2.6}.
\end{proof}

\begin{definition}\label{sga3:XXI:7.2.6}
Let \(\cR\) be any root datum.  The \emph{length} of
\(\alpha\in R\), denoted \(\operatorname{long}(\alpha)\), is the
length of \(\alpha\) in its irreducible component.
\end{definition}

\begin{lemma}\label{sga3:XXI:7.2.7}
There is a unique group homomorphism
\[
 u:\Gamma_0(R)\longrightarrow\Gamma_0(R^\ast)
\]
such that
\[
 u(\alpha)=\operatorname{long}(\alpha)\alpha^\ast
 \qquad(\alpha\in R).
\]
\end{lemma}

\begin{proof}
By \ref{sga3:XXI:3.5.5}, it is enough to verify that, whenever
\(\alpha,\beta,\alpha+\beta\in R\),
\[
 \operatorname{long}(\alpha)\alpha^\ast+
 \operatorname{long}(\beta)\beta^\ast
 =
 \operatorname{long}(\alpha+\beta)(\alpha+\beta)^\ast.
\]
But \(\alpha,\beta,\alpha+\beta\) belong to the same irreducible
component of \(R\), by \ref{sga3:XXI:7.1.2}; the assertion is
therefore reduced to \ref{sga3:XXI:1.2.2}.
\end{proof}

\begin{remark}\label{sga3:XXI:7.2.8}
Let \(u\) be as in \ref{sga3:XXI:7.2.7}.  For
\(\alpha,\beta\in R\),
\[
 (u(\alpha),\beta)=(u(\beta),\alpha).
\]
Indeed, this amounts to
\[
 \operatorname{long}(\alpha)(\alpha^\ast,\beta)
 =
 \operatorname{long}(\beta)(\beta^\ast,\alpha),
\]
which is clear when \(\alpha\) and \(\beta\) are orthogonal.  If
they are not orthogonal, they lie in the same irreducible component
of \(R\), and the assertion reduces to \ref{sga3:XXI:1.2.1},
formula~\SGARefLink{sga3:XXI:eq:9}{(9)}.
\end{remark}

\begin{remark}\label{sga3:XXI:7.2.9}
The symmetric bilinear form
\[
 (x,y)\longmapsto(u(x),y)
\]
is positive and nondegenerate on \(\Gamma_0(R)\).
\end{remark}

\begin{proof}
Let \(R_i\) be the irreducible components of \(R\).  Then
\[
 \Gamma_0(R)=\prod_i\Gamma_0(R_i),
\]
and the bilinear form \((u(x),y)\) is the product, on the groups
\(\Gamma_0(R_i)\), of the forms
\[
 2\ell(\alpha_i)^{-1}(p(x),y),
\]
where \(\ell(\alpha_i)\) is the minimum of \(\ell(\alpha)\) for
\(\alpha\in R_i\).  These symmetric bilinear forms are positive
and nondegenerate by \ref{sga3:XXI:1.2.6}.
\end{proof}
